\documentclass{article}
\usepackage{amsmath}
\usepackage{graphicx}
\usepackage{float}

\usepackage[utf8]{inputenc}     % pdfLaTeX
\usepackage[T1]{fontenc}
\usepackage[magyar]{babel}

\usepackage{xfrac}
\usepackage{nicefrac}

\title{Trónharc, trónviszály}
\author{Haydu Lőrinc, Lengyel Zoltán Péter, Rovenszky Robin}
\date{Legutóbb frissítve: 2026. Április 14.}

\newtheorem{remark}{Megjegyzés}

\begin{document}

\maketitle

\section{Kezdetek}
\begin{itemize}
    \item I. István halála előtt Vazult megvakíttatta, hogy ne legyen alkalmas a trónra
    \item Vazul fiait elűzték
    \item Orseolo Péter $1038-1041$. (István nővérét veszi feleségül, velencei dózse fia) 
    \begin{itemize}
        \item Németellenes
        \item Pogányüldöző
        \item Pécsi székesegyház
        \item István halála után Gizellát kisemmizi, megfosztja vagyonától és száműzi vidékre
        \item Idegeneket hív be királyi udvartartásra
        \item Aba Sámuel elűzi a trónról (István tesvérével házasodott)
    \end{itemize}
        
    \item Aba Sámuel $1041-1044$
    \begin{itemize}
        \item Német-Római császár és Orseolo Péter
        \item 1042. Osztrák Őrgrófságra betört $\rightarrow$ III: Henrik császár folyamatos támadásai
        \item Szent István politikájával szakít (nem üldözi a pogányokat)
        \item Balkáni bogumilok beengedése
        \item Parasztokkal tölti a szabadidejét
        \item Összeesküvések és III. Henrik császár, ménfői csata $\rightarrow$ elűzik, Orseolo visszakerül a trónra
    \end{itemize}
    
    \item Orseolo Péter (másodjára) $1044-1046$
    \begin{itemize}
        \item németek hűbérese
        \item bajor törvény 
        \item adók újbóli kivetése
        \item Politikája nem változott, nem tanult $\rightarrow$ három összeesküvés (Vata-féle pogánylázadás $\rightarrow$ kereszténypapok mészárlása 
        \begin{itemize}
            \item Gellért püspök (csanádi püspök) és Kelen-hegy
        \end{itemize}
        
        \item \textbf{Mártírhalál/Vértanú}: Meghal valamilyen fontos ügyért. A mártír feláldozza magát, a vértanú nem feltétlenül, csak fontos ügyért hal meg. Keresztény kifejezések. 
        
    \end{itemize}
    \item I. András $1046-1060$.
    \begin{itemize}
        \item \textbf{Dukátus} Bélával: Vagyis Béla megkapta az ország $\nicefrac{1\,}{\,3}\,$-át.
        \item Felesége Anasztázia (orosz támogatás)
        \item Eredetileg öccsét akarta örökösnek (Béla), de később fia született (Salamon).
        \item Korona vagy kard 
        \begin{itemize}
            \item Kard - lovagság
            \item Korona - királyság, DE! lemészárolta volna öccsét
        \end{itemize}
        \item Bélával való csatában meghal (Tihanyi apátságba temetik el) 
        \begin{itemize}
            \item Tihanyi apátság alapítólevele: 1055, legrégebbi magyar szórványemlék
        \end{itemize}
    \end{itemize}
    
    \item I. Béla $1060-1063$

    \item 11. évig nincs folytonos uralkodó
    
    \item I. Géza $1074-1077$
    \begin{enumerate}
        \item Salamon unokatestvére
        \item VIII. Gergely pápa nem támogatja $\rightarrow$ bizánci császár elismerése kell (Korona alsó része) 
    \end{enumerate}
    
    
\end{itemize}

\subsection{Szent korona}
\begin{itemize}
    \item Koronázási ékszer
    \item Hartvik legenda: II. Szilveszter pápától származik, I. Istvánt koronázták meg vele
    \item Két részből áll 
    \begin{itemize}
        \item Alsó abroncs - görög korona 
        \begin{itemize}
            \item Bizánc, női fejék
            \item Központi részén Krisztus, Bizánci császárok, Géza fejedelem, görög szentek
        \end{itemize}
        
        \item Felső keresztpántos rész - latin korona 
        \begin{itemize}
            \item Apostolok díszítik
        \end{itemize}
        
        \item Egyesítés feltehetően III. Béla idején (Istvánt csak a felsővel koronázták)
        \item Tetején kereszt, eredetileg egyenes de elferdült
    \end{itemize}
    
    \item 1440. Szilágyi Erzsébet szolgálójával raboltatja el $\rightarrow$ III. Frigyeshez kerül $\rightarrow$ Mátyásnak $80000$ arany forintba kerül visszavásárolniaű
    
    \item Volt Bécsben, Prágában, Pozsonyban, Budán, Orsovai mocsárban, II. vh alatt benzineshordóban Ausztriában, 1953 USA-ban kerül, 1978-ban kerül vissza hozzánk 
    
    \item 2000-ig a Nemzeti Múzeumban, azóta Parlamentben

    \item XVI. század: Szent Korona tan (A Szent Korona a királyi hatalmat testesíti meg)

    \item Utolsó koronázás 1916. IV. Károly, Ferenc József halála után

\end{itemize}

\subsection{Koronázási ékszerek}
\begin{itemize}
    \item Koronázási palást
    \item Országalma
    \item Koronázási kard
    \item Jogar
    \item Szent Korona (Hatalom megszerzése $\rightarrow$ hatalom)
\end{itemize}


\end{document}